% ===== 写作阶段开关 =====
\newif\ifoutlineguidance
\outlineguidancetrue  % 改为 \outlineguidancefalse 可隐藏“写作提示”。
\newif\ifdraftlayout
\draftlayouttrue      % 写作预览不插空白页；定稿前改为 false 并核对学院要求。
\newif\iffullfrontmatter
\fullfrontmatterfalse % 当前仅摘要/大纲，默认不输出正式授权声明和全套封面。

% 三级目录：章 -> 节 -> 小节，例如 4 / 4.3 / 4.3.2。
\setcounter{secnumdepth}{2}
\setcounter{tocdepth}{2}
\titlecontents{subsection}[2\ccwd]{\songti\zihao{-4}}
  {\thecontentslabel\hspace{0.5\ccwd}}{}
  {\hspace{.5em}\titlerule*{.}\contentspage}
\hypersetup{hidelinks,pdfdisplaydoctitle=true,
 pdftitle={异构CAD形变特征转换技术研究与模型转换实现：摘要与三级目录初稿},
 pdfauthor={作者待填写},pdfsubject={硕士论文写作框架}}
\emergencystretch=2em
\raggedbottom
% 非正式预览采用连续起章，保留模板的其余页边距、字体级别及页眉页脚。
\ifdraftlayout
 \makeatletter\@openrightfalse\makeatother
 \let\cleardoublepage\clearpage
 % 仅提纲预览压缩标题留白；关闭 draftlayout 后恢复原模板间距。
 \ctexset{section={beforeskip=10pt,afterskip=8pt},subsection={beforeskip=8pt,afterskip=6pt}}
\fi
% 这些提示是提纲说明，不是最终论文正文。
\newcommand{\chapterguide}[1]{%
 \ifoutlineguidance\par\noindent{\small\color{black!70}\textbf{本章写作重点：}#1}\par\medskip\fi}
\newcommand{\outlinehint}[1]{%
 \ifoutlineguidance\par\noindent{\small\color{black!70}\textbf{写作提示：}#1}\par\fi}
